\newpage
\begin{center}
{\fontsize{15pt}{18pt}\selectfont \textbf{XÂY DỰNG MÔ HÌNH KHUYẾN NGHỊ LAI KẾT HỢP NHÂN TỬ HOÁ MA TRẬN VÀ MẠNG LƯỚI THẦN KINH SÂU}\par}
\vspace{0.5cm}
{\fontsize{16pt}{19.2pt}\selectfont \textbf{TÓM TẮT}\par}
\end{center}
\vspace{0.2cm}

Nhân tử hóa ma trận (Matrix Factorization -- MF) mô hình hóa tương tác giữa khách hàng và sản phẩm bằng tích vô hướng của hai vector ẩn \parencite{koren2009matrix}: một phép kết hợp tuyến tính, và chỉ học được từ những sản phẩm đã có người mua. Với dữ liệu thời trang -- cực thưa và có nhiều sản phẩm mới ra mắt liên tục -- cả hai điều này đều là giới hạn. Dự án xây dựng \textbf{NeuMF-F}, một mô hình lai giữ cấu trúc hai nhánh của Neural Matrix Factorization \parencite{he2017ncf} -- nhánh tuyến tính Generalized Matrix Factorization (GMF) và nhánh phi tuyến Multi-Layer Perceptron (MLP), hợp nhất ở lớp dự đoán -- nhưng biểu diễn mỗi khách hàng và sản phẩm bằng Embedding mã ID cộng với phép chiếu của đặc trưng: thuộc tính và mô tả văn bản của sản phẩm, thông tin khách hàng, doanh số và thời gian. Khi huấn luyện, ID của sản phẩm được bỏ ngẫu nhiên để mô hình học cách chấm điểm cả những sản phẩm chưa từng được bán.

Thực nghiệm dùng bộ dữ liệu H\&M Personalized Fashion Recommendations \parencite{hm2022kaggle} (31,8 triệu giao dịch). Hai mẫu khách hàng không giao nhau được lấy tái lập được theo hàm băm: mẫu phát triển A (21.599 khách) chỉ dùng để tinh chỉnh trên tập xác thực, mẫu kiểm định B (21.030 khách) chỉ dùng một lần cho đánh giá cuối. Dữ liệu được chia theo thời gian; lọc $k$-core chỉ dùng dữ liệu trước mốc kiểm thử; tập ứng viên gồm cả sản phẩm mới (\VData{Btest}{newpct}\% sản phẩm đúng ở tập kiểm thử của mẫu B là sản phẩm chưa từng bán trước mốc); mọi đặc trưng được tính theo thời điểm, không dùng thông tin tương lai. Mười bốn mô hình -- NeuMF-F, từng nhánh GMF-F và MLP-F, hợp nhất muộn LateFusion-F, NeuMF chỉ dùng ID cùng hai nhánh của nó, và bảy baseline Random, Most Popular, MostPopular-Recent, gợi ý theo nội dung, ItemKNN, UserKNN, BPR-MF -- được tinh chỉnh với cùng ngân sách. Đánh giá cuối theo kế hoạch đăng ký trước, có khoá tập kiểm thử: Full Ranking, 5 seed, kiểm định Wilcoxon theo người dùng kèm khoảng tin cậy bootstrap và hiệu chỉnh Holm trên họ 10 so sánh.

Trên tập kiểm thử của mẫu B (\VData{Btest}{users} khách hàng), NeuMF-F đạt NDCG@10 = \VRes{NeuMFF}{NDCG10}, xếp \VRank{NeuMFF}{NDCG10} trong 14 mô hình; NeuMF chỉ dùng ID đạt \VRes{NeuMF}{NDCG10}, BPR-MF \VRes{BPR}{NDCG10}, UserKNN \VRes{UKNN}{NDCG10}, MostPopular-Recent \VRes{RecPop}{NDCG10}. \VSig{all}{summary} Trên nhóm sản phẩm mới, NeuMF-F đạt NDCG@10 = \VRes{NeuMFF}{NDCG10new}, trong khi mọi mô hình chỉ dùng ID bằng 0 theo cấu tạo. Giai đoạn phát triển trước đó (NeuMF chỉ dùng ID, đánh giá trên chính mẫu phát triển) được báo cáo như lịch sử phát triển: NeuMF không vượt MF đã tinh chỉnh -- lý do đề tài chuyển sang bổ sung thông tin thay vì chỉ tăng độ phức tạp của mô hình. Hệ thống demo suy diễn (FastAPI) minh họa và so sánh các mô hình trên từng khách hàng.
